\paragrafos{Παραδείγματα επίλυσης προβλημάτων: μάθε από τα λυμένα}{Α.4.3, «Παραδείγματα επίλυσης προβλημάτων»}

\begin{anagke}
Η πρώτη εβδομάδα στο γυμνάσιο τελειώνει και οι εννέα της παρέας, η Ζωή, ο Ηλίας, η Ρήνη, ο Στέλιος, η Φανή, ο Βασίλης, ο Νίκος, η Αγγελική και η Πολυξένη, αποφασίζουν να γιορτάσουν: κάθε παιδί θα δώσει το χέρι του σε όλα τα υπόλοιπα. «Πόσες χειραψίες θα γίνουν;» ρωτάει ο Ηλίας. Ο Βασίλης απαντάει με σιγουριά: «Εννιά επί εννιά, ογδόντα ένα». Η Πολυξένη τον κοιτάζει: «Και ο καθένας δίνει το χέρι του στον εαυτό του;».

Έτσι αρχίζει ένα πρόβλημα που δεν λύνεται με μία απλή εξίσωση, αλλά με σκέψη: μέτρηση, σχήμα, πίνακας, ανάποδη πορεία. Η Αγγελική σχεδιάζει εννιά τελείες σε κύκλο και ενώνει κάθε ζευγάρι. Ο Στέλιος προτείνει να μετρήσουν ένας-ένας. Και η Ρήνη τους προκαλεί: «Στο ταμείο θα πρέπει να βρούμε και πόσα χρήματα είχα, αφού ξόδεψα τα μισά για να αγοράσω τετράδια και κασετίνα».

Στην παράγραφο αυτή η παρέα λύνει παραδείγματα με διαφορετικά είδη εξισώσεων και τρόπων σκέψης. Η σιγουριά του Βασίλη θα ελεγχθεί με τον πιο απλό τρόπο: μετράμε.
\end{anagke}

\begin{lexi}{διαδοχικοί, προκαταβολή}{διαδοχικοί, προκαταβολή}
\lexiword{Διαδοχικοί} είναι οι αριθμοί που ακολουθούν ο ένας τον άλλον, χωρίς κενό: $\nu,\ \nu+1,\ \nu+2$. Η \lexiword{προκαταβολή} είναι χρήματα που δίνονται «προ-καταβάλλονται», δηλαδή πριν από την εργασία. Το \lexiword{μεροκάματο} είναι η αμοιβή για μία «μέρα» δουλειάς.
\end{lexi}

\begin{metafrasi}
Τρεις διαδοχικοί φυσικοί. & $x,\ x+1,\ x+2$\\
Η προκαταβολή είναι τα $\frac{2}{5}$ της αμοιβής. & $\dfrac{2}{5}\cdot(4x)=160$\\
Ξόδεψε τα μισά και κράτησε το υπόλοιπο. & $x-\dfrac{x}{2}=\dfrac{x}{2}$\\
Εννέα άνθρωποι, ο καθένας με τους άλλους. & $9\cdot8:2$\\
\end{metafrasi}

\tmima{Η θεωρία}

Στα προβλήματα που ακολουθούν δεν υπάρχει νέος κανόνας, αλλά καλό είναι να έχουμε ένα «ρεπερτόριο» τρόπων σκέψης:
\begin{itemize}[leftmargin=1.2em,itemsep=0.3ex,topsep=2pt]
\item \textbf{Μία ισότητα με λόγια πρώτα.} Γράφω «τα χρήματα που ξοδεύτηκαν = το κόστος των αγορών» και μετά αντικαθιστώ τα μέρη με τον άγνωστο.
\item \textbf{Ανάποδη πορεία.} Ξεκινάω από το τελικό αποτέλεσμα και αναιρώ τις πράξεις.
\item \textbf{Σχήμα ή πίνακας.} Βοηθάει όταν υπάρχουν πολλά ποσά, ή όταν θέλω να μετρήσω συνδυασμούς.
\item \textbf{Δύο τρόποι.} Όταν μπορώ, λύνω το πρόβλημα με δύο τρόπους και συγκρίνω.
\end{itemize}
Για το άθροισμα $1+2+3+\dots+(\nu-1)$ βοηθάει να ζευγαρώνουμε από τα άκρα: $(1+9)+(2+8)+(3+7)+(4+6)+5=45$. Για τις χειραψίες $\nu$ ατόμων, ο καθένας δίνει το χέρι σε $\nu-1$ άλλους, αλλά κάθε χειραψία μετρήθηκε δύο φορές, άρα ο συνολικός αριθμός είναι $\dfrac{\nu\cdot(\nu-1)}{2}$.

\begin{prosoxi}
\begin{itemize}[leftmargin=1.2em,itemsep=0.3ex,topsep=2pt]
\item Η χειραψία είναι μεταξύ \emph{δύο} ανθρώπων: την έχω μετρήσει δύο φορές αν κάνω $\nu\cdot(\nu-1)$. Διαιρώ με το $2$.
\item Όταν γράφω «το μισό» ή «τα $\frac{2}{5}$» πρέπει να ξέρω \emph{του τι}: του συνόλου ή του υπολοίπου.
\item Η εξίσωση επαληθεύεται, αλλά και το πρόβλημα: ο αριθμός των μαθητών δεν μπορεί να είναι $2{,}5$.
\end{itemize}
\end{prosoxi}

\tmima{Λυμένα παραδείγματα}

\paradeigma{Τα μισά χρήματα}
\tixero Η Ρήνη ξόδεψε τα μισά της χρήματα για $3$ τετράδια των $1{,}5$ € και μια κασετίνα των $4{,}5$ €.\par
\tithelo Πόσα χρήματα είχε αρχικά.\par
\possindeo Έστω $x$ τα χρήματά της. «Τα χρήματα που ξόδεψε = το κόστος των αγορών», δηλαδή $\dfrac{x}{2}=3\cdot1{,}5+4{,}5=9$. Άρα $x=9\cdot2=18$.\par
\apantisi $18$ €.\par
\elegxos Τα μισά των $18$ είναι $9$ και $3\cdot1{,}5+4{,}5=9$. Σωστό.

\paradeigma{Η δεξαμενή του χωριού}
\tixero Η δεξαμενή χωράει $5\,400$ λίτρα. Καταναλώνονται $300$ λίτρα την ημέρα. Έξι ημέρες μετά περιορίζουν την κατανάλωση σε $200$ λίτρα την ημέρα.\par
\tithelo Πόσες ακόμη ημέρες θα φτάσει το νερό.\par
\possindeo Έστω $x$ οι επιπλέον ημέρες. «Νερό που καταναλώθηκε = νερό δεξαμενής»: $300\cdot6+200\cdot x=5\,400$. Άρα $1\,800+200x=5\,400$, $200x=3\,600$, $x=18$.\par
\apantisi $18$ ακόμη ημέρες.\par
\elegxos $300\cdot6+200\cdot18=1\,800+3\,600=5\,400$.

\paradeigma{Το μεροκάματο και η προκαταβολή}
\tixero Η θεία Μαρία δουλεύει $4$ ημέρες και παίρνει προκαταβολή τα $\frac{2}{5}$ της αμοιβής της, δηλαδή $160$ €.\par
\tithelo Το μεροκάματο.\par
\possindeo Έστω $x$ το μεροκάματο. Η αμοιβή είναι $4x$ και η προκαταβολή $\frac{2}{5}\cdot4x=160$. Άρα $4x=160:\frac{2}{5}=160\cdot\frac{5}{2}=400$ και $x=100$.\par
\apantisi $100$ € την ημέρα.\par
\elegxos Αμοιβή $4\cdot100=400$ και τα $\frac{2}{5}$ του $400$ είναι $160$.

\paradeigma{Οι χειραψίες της παρέας}
\tixero Εννέα παιδιά δίνουν το χέρι το καθένα με όλα τα άλλα.\par
\tithelo Πόσες χειραψίες γίνονται.\par
\possindeo \emph{1ος τρόπος:} Το πρώτο παιδί δίνει το χέρι σε $8$, το δεύτερο σε $7$ ακόμη (όχι στο πρώτο, το έχει κάνει), \dots, το όγδοο σε $1$ και το ένατο σε κανένα. Άθροισμα: $8+7+6+5+4+3+2+1=36$. \emph{2ος τρόπος:} Κάθε παιδί δίνει το χέρι σε $8$ άλλα, άρα $9\cdot8=72$, αλλά κάθε χειραψία μετρήθηκε δύο φορές, άρα $72:2=36$.\par
\apantisi $36$ χειραψίες.\par
\elegxos Και οι δύο τρόποι δίνουν $36$. Το $81$ του Βασίλη ήταν λάθος: μέτρησε και τον εαυτό.

\paradeigma{Διαδοχικοί αριθμοί}
\tixero Το άθροισμα τριών διαδοχικών φυσικών είναι $78$.\par
\tithelo Οι αριθμοί.\par
\possindeo Αν ο μικρότερος είναι $x$, οι άλλοι είναι $x+1$ και $x+2$. $x+(x+1)+(x+2)=78$, δηλαδή $3x+3=78$, $3x=75$, $x=25$.\par
\apantisi $25,\,26,\,27$.\par
\elegxos $25+26+27=78$.

\tmima{Ασκήσεις}

\zoni{Το κατάλαβα}{διαδοχικοί, κλάσματα, ανάποδη πορεία}

\ex Βρες τρεις διαδοχικούς φυσικούς με άθροισμα $111$.

\ex Βρες δύο αριθμούς με άθροισμα $51$, αν ο ένας είναι διπλάσιος του άλλου.

\ex Βρες τέσσερις διαδοχικούς άρτιους αριθμούς με άθροισμα $100$.

\ex Τα $\frac{3}{5}$ ενός αριθμού είναι $36$. Ποιος είναι ο αριθμός;

\ex Ο Ηλίας ξόδεψε το $\frac{1}{3}$ των χρημάτων του και ακόμη $4$ € και του έμειναν $10$ €. Πόσα είχε αρχικά;

\zoni{Το χρησιμοποιώ}{ηλικίες, δεξαμενές, αγώνες}

\ex \textbf{Τα αδέλφια.} Η Πολυξένη είναι $5$ χρόνια μικρότερη από τον αδελφό της και μαζί έχουν $29$ χρόνια. Πόσο χρονών είναι η καθεμία;

\ex \textbf{Η δεξαμενή του σχολείου.} Η δεξαμενή χωράει $4\,800$ λίτρα. Καταναλώνονται $250$ λίτρα την ημέρα και μετά από $8$ ημέρες η κατανάλωση πέφτει στα $100$ λίτρα την ημέρα. Πόσες ακόμη ημέρες θα φτάσει το νερό;

\ex \textbf{Το πρωτάθλημα.} Έξι ομάδες παίζουν μία φορά η καθεμία με όλες τις άλλες. Πόσοι αγώνες γίνονται; Και αν ήταν δέκα ομάδες;

\ex \textbf{Το λάδι.} Τα $\frac{3}{8}$ του μπουκαλιού του παππού είναι $1{,}2$ λίτρα. Πόσα λίτρα χωράει το μπουκάλι;

\ex \textbf{Η προκαταβολή.} Η θεία Μαρία δούλεψε $6$ ημέρες και πήρε προκαταβολή το μισό της αμοιβής της, $270$ €. Πόσο ήταν το μεροκάματο;

\ex \textbf{Το ταξί.} Μια διαδρομή με ταξί κοστίζει $3{,}5$ € πάγιο και $0{,}8$ € το χιλιόμετρο. Η Αγγελική πλήρωσε $13{,}1$ €. Πόσα χιλιόμετρα έκανε;

\ex \textbf{Οι δρομείς.} Δύο δρομείς ξεκινούν τον ίδιο χρόνο από τα δύο άκρα μιας διαδρομής $5$ χλμ. και τρέχουν ο ένας προς τον άλλον με $8$ και $12$ χλμ. την ώρα. Σε πόση ώρα θα συναντηθούν;

\zoni{Το σκέφτομαι}{όταν ο πρώτος τρόπος δεν φτάνει}

\ex Ο Βασίλης λέει ότι σε μια ομάδα $12$ ατόμων γίνονται $12\cdot11=132$ χειραψίες. Πού κάνει λάθος; Πόσες είναι σωστά;

\ex Σκέφτηκα έναν αριθμό. Του πρόσθεσα το μισό του και βρήκα $36$. Ποιος ήταν;

\ex Πόσα ζευγάρια $(x,y)$ φυσικών αριθμών έχουν άθροισμα $10$; Γιατί το πρόβλημα έχει πολλές λύσεις;

\ex Σε ένα τουρνουά με διπλό γύρο (η κάθε ομάδα παίζει δύο φορές με την άλλη, μία εντός και μία εκτός) συμμετέχουν $8$ ομάδες. Πόσοι αγώνες γίνονται;

\ex Υπολόγισε το $1+2+3+\dots+20$ ζευγαρώνοντας τους όρους από τα άκρα.

\zoni{Για το σπίτι}{ένα δικό σου πρόβλημα}

\ex Φτιάξε ένα πρόβλημα με τρεις διαδοχικούς αριθμούς, π.χ. τρία γειτονικά σπίτια της οδού, και λύσε το.

\ex Ψάξε στο διαδίκτυο ποιος ήταν ο Διόφαντος και γράψε μια πρόταση για το πώς ονόμαζε τον άγνωστο.
